% !TeX encoding = UTF-8
\SetChapterInfo{subtitle={从极大理想到谱积分与调和分析},keywords={Gelfand 表示; 谱测度; 自伴算子; Friedrichs; 紧群; Tannaka; Wiener},prereq={预解式; Cayley 变换; 紧算子; Riesz--Markov; Stone--Weierstrass},goals={构造谱表示; 控制无界演算定义域; 连接群表示与 Fourier 分析}}
\chapter[赋范环和谱表示]{赋范环\\和谱表示}\label{ch:yosida-11}
\ChapterSource{按 \cite{yosida-zh,yosida-en} 第十一章十六节组织，2026 年 10 月 27 日补写。直接查看英文印刷 294--362 页扫描图，第 362 页上部为本章结尾、下部开始第十二章，故更正本章英文范围。中译本参考范围为 307--376 页，选页及证明边界详见本章 COVERAGE。}

第十章的紧性使非零谱问题归结为有限维空间，本章转向一般正规与自伴算子。第一条线索是把交换 Banach 代数的元素表示为函数；第二条线索是用投影值测度把这些函数重新解释为算子；第三条线索则把谱分解用于群表示、正定函数和 Tauber 型逼近。

本章 $H$ 为非零复 Hilbert 空间，内积对第一变量线性。$\mathcal A$ 表示代数，$A$ 或 $S$ 表示算子，$E$ 表示投影值测度。实轴谱族统一采用右连续约定 $E_t=E((-\infty,t])$。Fourier 变换仍取 $\widehat f(\xi)=\int e^{-ix\xi}f(x)\,\dd x$，原书约定在对应公式中换算。测度表示的外部前提在使用处注明。

\Needspace{10\baselineskip}
\section{赋范环的极大理想}\label{sec:yosida-11-01}
\noindent\textit{原书起始页：中译本 309；英文版 295。}

\begin{definition}[Banach 代数与特征]
复 Banach 空间 $\mathcal A$ 若具有结合、双线性的乘法且 $\|ab\|\le\|a\|\|b\|$，称为 Banach 代数。以下先假定其交换、有单位 $1$ 且 $\|1\|=1$。特征指非零复线性乘法泛函 $\chi:\mathcal A\to\mathbb C$，全体记为 $\Delta(\mathcal A)$。真理想中的极大元称为极大理想。
\end{definition}

本节的理想是满足 $\mathcal A M\subset M$ 的线性子空间 $M$；“真”表示 $M\ne\mathcal A$。有单位时，这等价于 $1\notin M$。

\begin{theorem}[极大理想与特征对应]\label{thm:y11-characters}
每个真理想包含于某极大理想。每个极大理想 $M$ 闭，$\mathcal A/M\cong\mathbb C$，且唯一对应 $\chi\in\Delta(\mathcal A)$ 使 $M=\ker\chi$。每个特征自动连续，$\|\chi\|=\chi(1)=1$。
\end{theorem}
\begin{proof}
用包含给定真理想的真理想族及 Zorn 引理取极大元；链的并仍真，因为它不含 $1$。若 $\overline M=\mathcal A$，则有 $a\in M$ 使 $\|1-a\|<1$；Neumann 级数使 $a$ 可逆，导致 $1\in M$，矛盾。闭包为理想，由极大性 $M=\overline M$。

商代数是复 Banach 除代数。对其元素 $b$，谱非空：若预解式处处存在，则在无穷远趋零，逐连续泛函应用 Liouville 定理会使预解式恒为零，矛盾。取 $\lambda\in\sigma(b)$，除代数中不可逆元素只能为零，故 $b=\lambda1$。这给出 Gelfand--Mazur 结论及商特征。反之，特征的核的商是一维域，所以核极大。

对任意特征，$\chi(1)=1$；$a-\chi(a)1$ 不可逆，否则对其逆应用 $\chi$ 得矛盾。又 $|\lambda|>\|a\|$ 时 Neumann 级数给出逆，故 $|\chi(a)|\le\|a\|$。结合单位值，得到范数为 $1$，无需预先假定连续性。
\end{proof}

\begin{example}[连续函数代数]
若 $K$ 紧 Hausdorff，则 $C(K)$ 的特征恰为点值 $f\mapsto f(t)$。若极大理想 $M$ 的函数无公共零点，紧性给出 $f_1,\ldots,f_n\in M$ 使 $h=\sum|f_j|^2>0$；$h\in M$ 且可逆，矛盾。因此 $M\subset\{f:f(t)=0\}$，极大性使它们相等。
\end{example}

\begin{exercise}[单位化]\label{ex:y11-unitization}
对无单位 Banach 代数 $\mathcal B$，在 $\mathbb C\oplus\mathcal B$ 上写出乘法与范数，使 $(1,0)$ 成为单位，并验证乘法估计。
\end{exercise}
\begin{solution}
定义 $(\lambda,a)(\mu,b)=(\lambda\mu,\lambda b+\mu a+ab)$，范数 $|\lambda|+\|a\|$。三角不等式和 $\|ab\|\le\|a\|\|b\|$ 给出积的范数至多两范数之积；完备性来自直和完备性。
\end{solution}

\begin{exercise}[交换性的位置]\label{ex:y11-noncommutative}
说明矩阵代数 $M_2(\mathbb C)$ 不能套用本节极大理想到复数特征的对应。
\end{exercise}
\begin{solution}
$M_2(\mathbb C)$ 的双边理想只有零与全代数，故零理想极大，但商仍为矩阵代数而非域。若有复特征，其核只能为零，得到从四维复空间到一维空间的单射，矛盾。交换性用于将极大理想商变成除代数。
\end{solution}

\Needspace{10\baselineskip}
\section{根、半单性}\label{sec:yosida-11-02}
\noindent\textit{原书起始页：中译本 312；英文版 298。}

\begin{definition}[Gelfand 变换与根]
对 $a\in\mathcal A$ 定义 $\widehat a(\chi)=\chi(a)$。赋 $\Delta(\mathcal A)$ 逐点收敛拓扑，称为特征空间。根为 $\operatorname{rad}\mathcal A=\bigcap_{\chi}\ker\chi$；根为零时称半单。
\end{definition}

\begin{theorem}[交换 Banach 代数的函数表示]\label{thm:y11-gelfand}
$\Delta(\mathcal A)$ 为紧 Hausdorff 空间；Gelfand 变换是到 $C(\Delta(\mathcal A))$ 的连续代数同态，并且
\[
 \sigma_{\mathcal A}(a)=\{\chi(a):\chi\in\Delta(\mathcal A)\},\qquad
 \|\widehat a\|_\infty=r(a)=\lim_n\|a^n\|^{1/n}.
\]
其核为根。因此交换代数中，根由所有拟幂零元素组成；半单性等价于 Gelfand 变换单射，不自动等价于等距。
\end{theorem}
\begin{proof}
把特征嵌入 $\prod_{a\in\mathcal A}\{z:|z|\le\|a\|\}$。线性、乘法与 $\chi(1)=1$ 均为闭条件，故特征空间是紧乘积的闭子集，且坐标函数分离点。若 $a-\lambda1$ 不可逆，其生成理想真，故包含于某极大理想；相应特征满足 $\chi(a)=\lambda$。反向已在上一节证明。谱半径公式的证明与第八章相同：Neumann 展开与预解式的 Cauchy 系数估计只用代数乘法及完备性，因此适用于 $\mathcal A$。由谱等式即得范数等式，核和根的断言随定义得到。
\end{proof}

\begin{theorem}[交换算子 C 星代数]\label{thm:y11-commutative-cstar}
设 $\mathcal A\subset\mathcal L(H)$ 为含单位、范数闭、对伴随封闭的交换子代数。则 Gelfand 变换是到 $C(\Delta(\mathcal A))$ 的等距满 $*$ 同构。
\end{theorem}
\begin{proof}
每个元素正规。由 $\|T^*T\|=\|T\|^2$ 及交换性，反复平方得 $\|T^{2^n}\|=\|T\|^{2^n}$，故 $r(T)=\|T\|$。于是 Gelfand 变换等距，像闭。若 $S=S^*$，$e^{itS}$ 酉，所以 $|e^{it\chi(S)}|\le1$ 对正负 $t$ 均成立，迫使 $\chi(S)$ 实。分解 $T$ 为实部与虚部，得 $\chi(T^*)=\overline{\chi(T)}$。像因而含常数、分离点且对共轭封闭，Stone--Weierstrass 定理使其一致稠密；闭性给出满射。
\end{proof}

\begin{exercise}[非零根]\label{ex:y11-radical}
在代数 $\mathbb C[\varepsilon]/(\varepsilon^2)$ 上取范数 $\|a+b\varepsilon\|=|a|+|b|$。求特征、根和谱。
\end{exercise}
\begin{solution}
特征使 $\chi(\varepsilon)^2=0$，故唯一特征为 $a+b\varepsilon\mapsto a$。根为 $\mathbb C\varepsilon$，$\sigma(a+b\varepsilon)=\{a\}$。非零元素 $\varepsilon$ 幂零，说明一般交换 Banach 代数不一定半单。
\end{solution}

\begin{exercise}[Wiener 逆定理的离散版本]\label{ex:y11-wiener-discrete}
证明 $\ell^1(\mathbb Z)$ 卷积代数的特征为 $a\mapsto\sum_n a_nz^n$，$|z|=1$；从而绝对收敛 Fourier 级数处处非零时，其倒数也有绝对收敛 Fourier 级数。
\end{exercise}
\begin{solution}
令 $z=\chi(\delta_1)$，因 $\delta_1\delta_{-1}=\delta_0$，有 $|z|\le1$ 且 $|z^{-1}|\le1$，故 $|z|=1$。由有限支集序列稠密及特征连续性得到级数公式。反向这些公式显然为特征。处处非零即在每个极大理想商中非零，所以元素可逆；其逆属于 $\ell^1$，给出所求倒数级数。
\end{solution}

\Needspace{10\baselineskip}
\section{有界正规算子的谱分解}\label{sec:yosida-11-03}
\noindent\textit{原书起始页：中译本 315；英文版 302。}

\begin{theorem}[连续函数演算]\label{thm:y11-continuous-calculus}
若 $N\in\mathcal L(H)$ 正规，则存在唯一等距含单位 $*$ 同态
\[
 C(\sigma(N))\longrightarrow\mathcal L(H),\qquad f\longmapsto f(N),
\]
把坐标函数 $z$ 送到 $N$，并满足 $\|f(N)\|=\|f\|_\infty$。
\end{theorem}
\begin{proof}
令 $\mathcal A=C^*(I,N)$，即 $N,N^*$ 的多项式的范数闭包。它交换，故应用定理 \ref{thm:y11-commutative-cstar}。映射 $\chi\mapsto\chi(N)$ 单射，因为这些多项式稠密。还需验证 $\sigma_{\mathcal A}(N)=\sigma_{\mathcal L(H)}(N)$。

若 $T\in\mathcal A$ 在 $\mathcal L(H)$ 中可逆，则 $S=T^*T$ 有 $cI\le S\le MI$，$c>0$。对自伴算子的二次型范数公式给出 $\|I-S/M\|\le1-c/M<1$，故 $S^{-1}$ 的 Neumann 级数属于 $\mathcal A$。于是 $T^{-1}=S^{-1}T^*\in\mathcal A$。取 $T=N-\lambda I$ 便得谱相等，故特征空间与 $\sigma(N)$ 同胚。逆 Gelfand 变换给出演算，唯一性由 $z,\bar z$ 多项式在 $C(\sigma(N))$ 中稠密得到。
\end{proof}

\begin{theorem}[正规算子的投影值测度]\label{thm:y11-normal-pvm}
存在唯一 Borel 投影值测度 $E$，使 $E(\sigma(N))=I$，
\[
 E(B)E(C)=E(B\cap C),\qquad
 N=\int_{\sigma(N)}z\,\dd E(z).
\]
$E$ 的可列可加性采用强算子拓扑。对任意有界 Borel 函数 $f$，
\[
 \|f(N)x\|^2=\int|f|^2\,\dd\mu_x,\qquad
 \mu_x(B)=(E(B)x,x).
\]
\end{theorem}
\begin{proof}
先固定 $x$。泛函 $f\mapsto(f(N)x,x)$ 为正，因为非负连续函数有连续平方根。用外部前提 Riesz--Markov 定理 \ref{thm:y4-riesz-markov} 得到有限正测度 $\mu_x$。映射 $f\mapsto f(N)x$ 满足
\[
 (f(N)x,g(N)x)=\int f\overline g\,\dd\mu_x.
\]
连续函数在 $L^2(\mu_x)$ 中稠密，因此该映射延为 $L^2(\mu_x)$ 到循环空间 $H_x=\overline{C(\sigma(N))x}$ 的酉映射。在此映射下 $N$、$N^*$ 分别为乘 $z$、乘 $\bar z$，所以 $H_x$ 为约化子空间。

以 Zorn 引理取极大正交循环空间族；若其正交和不为 $H$，补空间仍约化，可再取一个非零循环空间，矛盾。因此 $H=\bigoplus_\alpha H_{x_\alpha}$。在每个 $L^2(\mu_{x_\alpha})$ 上以乘 $1_B$ 定义 $E(B)$，再作正交直和。乘法关系和强可加性由指示函数及控制收敛得到。即使直和不可数，每个向量也只有至多可数个非零分量，故此步骤无可分性要求。相同构造定义 $f(N)$ 并给出范数公式。

若另有这样的投影值测度，积分多项式先给出同一连续演算。各向量的标量测度由连续积分唯一决定；极化后得到所有 $(E(B)x,y)$，故投影测度唯一。其支集恰为 $\sigma(N)$，否则一个非零连续函数被演算为零，会违背等距性。
\end{proof}

本节将原书的交换正规族结论完整证明到单个正规算子的演算及谱测度。有限个两两交换且连同伴随一起交换的正规算子也可对其生成的交换 C 星代数作同样构造，得到联合谱上的投影值测度；任意族的原书双交换子与范畴集构造不在此重写。

\begin{exercise}[正规性不可省略]\label{ex:y11-jordan-normal}
证明 $J=\left(\begin{smallmatrix}0&1\\0&0\end{smallmatrix}\right)$ 没有使 $J=\int z\,\dd E(z)$ 的正交投影值谱表示。
\end{exercise}
\begin{solution}
投影值积分与其伴随分别为乘 $z$ 与乘 $\bar z$，必交换，故必正规。但 $JJ^*=\operatorname{diag}(1,0)$，$J^*J=\operatorname{diag}(0,1)$，不相等。
\end{solution}

\begin{exercise}[投影的谱测度]\label{ex:y11-projection-pvm}
设 $P=P^*=P^2$。写出其谱测度和 $f(P)$。
\end{exercise}
\begin{solution}
$E(B)=1_B(0)(I-P)+1_B(1)P$，故 $f(P)=f(0)(I-P)+f(1)P$。若某个子空间为零，相应谱点可从支集中删去。
\end{solution}

\Needspace{10\baselineskip}
\section{酉算子的谱分解}\label{sec:yosida-11-04}
\noindent\textit{原书起始页：中译本 319；英文版 306。}

\begin{theorem}[单位圆上的谱积分]\label{thm:y11-unitary-pvm}
若 $U$ 酉，则其谱测度支于单位圆 $\mathbb T$，且对所有 $n\in\mathbb Z$，
\[
 U^n=\int_{\mathbb T}z^n\,\dd E(z),\qquad U^*=\int_{\mathbb T}\bar z\,\dd E(z).
\]
反之，单位圆上总质量为 $I$ 的投影值测度的坐标积分是酉算子。
\end{theorem}
\begin{proof}
$U$ 正规，$U^*U=I$ 使连续演算中 $|z|^2=1$ 于谱上；也可对 $|\lambda|>1$ 用 $U$ 的 Neumann 级数，对 $|\lambda|<1$ 用 $U^{-1}$ 的级数得谱包含。各幂由演算乘法得到，反向由积分 $z\bar z=1$ 得 $U^*U=UU^*=I$。
\end{proof}

选定 Borel 辐角 $\operatorname{Arg}:\mathbb T\to[0,2\pi)$，定义 $B=\operatorname{Arg}(U)$。则 $B$ 为有界自伴算子，$U=e^{iB}$，但这个对数依赖所选分支，不能据此宣称唯一。若用实参数谱族 $F_t=E(\{z:\operatorname{Arg}z\le t\})$，应把 $z=1$ 的质量计入 $t=0$；原书使用另一端点记法时须整体换算，不能重复计算该原子。

\begin{example}[平移的谱表示]
在 $\ell^2(\mathbb Z)$ 上取双边移位 $(Sx)_n=x_{n-1}$。圆周取归一化测度 $\dd\theta/(2\pi)$，则酉 Fourier 级数映射 $\mathcal Fx(\theta)=\sum_nx_ne^{in\theta}$ 将它变为乘 $e^{i\theta}$。因此谱为 $\mathbb T$，谱投影为 $\mathcal F^{-1}1_B(e^{i\theta})\mathcal F$。单点测度为零，所以没有特征值。
\end{example}

\begin{exercise}[谱投影的原子]\label{ex:y11-unitary-atom}
证明 $E(\{z_0\})H=\ker(U-z_0I)$。
\end{exercise}
\begin{solution}
由 $\|(U-z_0)x\|^2=\int|z-z_0|^2\,\dd\mu_x$，该范数为零当且仅当 $\mu_x$ 集中于 $z_0$。这又等价于 $E(\{z_0\})x=x$。
\end{solution}

\begin{exercise}[酉对数不唯一]\label{ex:y11-unitary-log}
给出两个不同的有界自伴算子 $B_1,B_2$，使 $e^{iB_1}=e^{iB_2}=I$。
\end{exercise}
\begin{solution}
取 $B_1=0$，$B_2=2\pi P$，其中 $P$ 为任意非零正交投影。连续演算给出 $e^{2\pi iP}=(I-P)+P=I$。
\end{solution}

\Needspace{10\baselineskip}
\section{单位分解}\label{sec:yosida-11-05}
\noindent\textit{原书起始页：中译本 322；英文版 309。}

\begin{definition}[实轴单位分解]
一族正交投影 $E_t$ 称为右连续单位分解，若
\[
 E_sE_t=E_{\min(s,t)},\quad
 \mathop{\mathrm{s\!-!lim}}_{t\to-\infty}E_t=0,\quad
 \mathop{\mathrm{s\!-!lim}}_{t\to+\infty}E_t=I,\quad
 \mathop{\mathrm{s\!-!lim}}_{h\downarrow0}E_{t+h}=E_t.
\]
它对应 $E((a,b])=E_b-E_a$ 的唯一投影值测度；反过来每个实轴投影值概率测度给出此谱族。这里的“单位分解”是 resolution of identity，与光滑单位分解不是同一概念。
\end{definition}

\begin{proposition}[从谱族到测度]
对 $x,y\in H$，$t\mapsto(E_tx,y)$ 有界变差，其全变差至多 $\|x\|\|y\|$。相应标量 Stieltjes 测度 $\mu_{x,y}$ 满足 $\mu_x=\mu_{x,x}\ge0$，并定义上述唯一投影值测度。
\end{proposition}
\begin{proof}
对递增分割 $t_j$，令 $P_j=E_{t_j}-E_{t_{j-1}}$，这些投影两两正交。因此
\[
 \sum_j|(P_jx,y)|=\sum_j|(P_jx,P_jy)|
 \le\Bigl(\sum_j\|P_jx\|^2\Bigr)^{1/2}
     \Bigl(\sum_j\|P_jy\|^2\Bigr)^{1/2}\le\|x\|\|y\|.
\]
标量 Stieltjes 延拓和极化得到所有 $\mu_{x,y}$。对每个 Borel 集 $B$，由此界及 Riesz 表示定义有界算子 $E(B)$，使 $(E(B)x,y)=\mu_{x,y}(B)$。区间上的投影乘法关系，依次固定一个变量再用有界单调类延拓到所有 Borel 集，得到 $E(B)E(C)=E(B\cap C)$ 及自伴性。标量可加性先给弱可加性；不交投影的正交性使余项的范数平方等于相应正测度余项，故为强可加性。
\end{proof}

\begin{theorem}[谱积分的定义域]\label{thm:y11-spectral-integral}
对有限值 Borel 函数 $f:\R\to\mathbb C$，定义
\[
 D(T_f)=\{x:\int|f|^2\,\dd\mu_x<\infty\},\qquad
 T_fx=\lim_{n\to\infty}\int f1_{\{|f|\le n\}}\,\dd E\,x.
\]
此极限存在，定义稠密闭算子，且 $\|T_fx\|^2=\int|f|^2\,\dd\mu_x$，$T_f^*=T_{\bar f}$。实值 $f$ 给出自伴算子。
\end{theorem}
\begin{proof}
对简单函数由正交性直接得到范数公式。用有界简单函数一致逼近有界 Borel 函数，公式传到有界演算。截断之差的平方范数恰为 $|f|^2$ 在相应层集的积分，故其 Cauchy 性等价于定义域条件。投影 $P_n=E(\{|f|\le n\})$ 强趋于 $I$，且 $P_nH\subset D(T_f)$，故定义域稠密。

若 $x_j\to x$、$T_fx_j\to y$，则 $f1_{\{|f|\le n\}}(E)x=P_ny$。于是截断平方积分有界于 $\|y\|^2$，单调收敛使 $x\in D(T_f)$，再令 $n\to\infty$ 得 $T_fx=y$。这证明闭性。显然 $T_{\bar f}\subset T_f^*$。若 $y\in D(T_f^*)$，在 $P_nH$ 上比较伴随得 $T_{\bar f}P_ny=P_nT_f^*y$，同样的积分估计使 $y\in D(T_{\bar f})$。故伴随恰相等。
\end{proof}

\begin{exercise}[强收敛与范数收敛]\label{ex:y11-projection-limits}
在 $\ell^2$ 上令 $P_n$ 为前 $n$ 个坐标的投影。证明 $P_n\to I$ 强收敛，但不按算子范数收敛。
\end{exercise}
\begin{solution}
$\|(I-P_n)x\|^2=\sum_{j>n}|x_j|^2\to0$，而 $\|(I-P_n)e_{n+1}\|=1$，所以 $\|I-P_n\|=1$。
\end{solution}

\begin{exercise}[无界积分]\label{ex:y11-integral-domain}
在 $L^2(\R)$ 的坐标谱测度上取 $f(t)=t^2$。求 $D(T_f)$，并说明 $D(T_f)$ 与 $D(T_t)$ 的关系。
\end{exercise}
\begin{solution}
$D(T_f)=\{u:\int t^4|u(t)|^2\,\dd t<\infty\}$，而 $D(T_t)$ 只要求二阶加权平方积分有限。前者真包含于后者；例如 $u(t)=1_{(1,\infty)}t^{-2}$ 属于后者而不属于前者。
\end{solution}

\Needspace{10\baselineskip}
\section{自伴算子的谱分解}\label{sec:yosida-11-06}
\noindent\textit{原书起始页：中译本 327；英文版 313。}

\begin{theorem}[无界自伴谱定理]\label{thm:y11-self-adjoint}
每个自伴算子 $A$ 都有唯一的实轴投影值测度 $E$，使
\[
 A=\int_\R t\,\dd E(t),\qquad
 D(A)=\{x:\int_\R t^2\,\dd\mu_x(t)<\infty\}.
\]
对于 $z\notin\R$，
\[
 (A-zI)^{-1}=\int_\R(t-z)^{-1}\,\dd E(t).
\]
\end{theorem}
\begin{proof}
由定理 \ref{thm:y7-cayley}，$U=(A-iI)(A+iI)^{-1}=I-2i(A+iI)^{-1}$ 酉，且 $\ker(I-U)=0$。取其单位圆谱测度 $P$，则 $P(\{1\})=0$。以下映射为 Borel 同构：
\[
 c(t)=\frac{t-i}{t+i}:\R\longrightarrow\mathbb T\setminus\{1\},\qquad
 c^{-1}(z)=i\frac{1+z}{1-z}
\]
\par\Needspace{6\baselineskip}
定义 $E(B)=P(c(B))$，按上一节构造自伴算子 $A_0=\int t\,\dd E(t)$，则
\[
 (A_0+iI)^{-1}=\frac{I-U}{2i}=(A+iI)^{-1}.
\]
相同逆算子的值域给出相同定义域，再取逆得 $A_0=A$。一般非实预解式由有界乘子 $(t-z)^{-1}$ 与 $t-z$ 相乘检验双侧逆。若另有谱测度表示 $A$，经 $c$ 推回后表示同一 $U$；酉谱测度唯一，故 $E$ 唯一。
\end{proof}

若 $z=e^{i\theta}$，直接计算得 $c^{-1}(z)=-\cot(\theta/2)$。中译本此处有半角校注；不能写成 $-\cot\theta$。端点 $z=1$ 的投影质量为零，而不要求 $1$ 不属于 $U$ 的谱：无界 $A$ 常使它成为连续谱点。

\begin{example}[动量算子的定义域]
令 $\mathcal F_0=(2\pi)^{-1/2}\mathcal F$ 为 $L^2(\R)$ 上的酉 Fourier 变换。对 $P=-i\,\dd/\dd x$，$D(P)=H^1(\R)$，有
\[
 \mathcal F_0P\mathcal F_0^{-1}=M_\xi,\qquad
 E_P(B)=\mathcal F_0^{-1}1_B(\xi)\mathcal F_0.
\]
先对 Schwartz 函数分部积分，再用 Fourier 加权范数闭包得到定义域相等。$C_c^\infty$ 上同一微分表达式只是其稠密限制，不能直接把该限制称作自伴算子。
\end{example}

\begin{exercise}[Cayley 的端点]\label{ex:y11-cayley-endpoint}
对 $A=M_t$ 作用于 $L^2(\R)$，证明 $1\in\sigma(c(A))$，但 $P(\{1\})=0$。
\end{exercise}
\begin{solution}
$c(t)$ 对有限实数从不等于 $1$，故单点投影为零；但 $c(t)\to1$ 当 $|t|\to\infty$。取支于 $(n,n+1)$ 的单位向量 $u_n$，则 $\|(c(A)-I)u_n\|\to0$，故不存在有界逆，$1$ 属于谱。
\end{solution}

\begin{exercise}[谱截断是图范数逼近]\label{ex:y11-spectral-core}
设 $P_n=E([-n,n])$。证明 $P_nx\to x$ 对每个 $x\in H$ 成立，且对 $x\in D(A)$ 有 $AP_nx\to Ax$。
\end{exercise}
\begin{solution}
两项误差的平方分别为 $\mu_x(\{|t|>n\})$ 与 $\int_{|t|>n}t^2\,\dd\mu_x$，均趋于零。因此 $\bigcup_nP_nH$ 在 $D(A)$ 的图范数中稠密。
\end{solution}

\Needspace{10\baselineskip}
\section{实算子和半有界算子、Friedrichs 定理}\label{sec:yosida-11-07}
\noindent\textit{原书起始页：中译本 329；英文版 316。}

\begin{definition}[实性与半有界性]
共轭算子 $C$ 是满足 $C^2=I$ 的反线性等距映射，例如 $L^2$ 上的复共轭。若 $CD(S)=D(S)$ 且 $SC=CS$，称 $S$ 对 $C$ 为实算子。稠密定义的对称算子 $S$ 称下半有界，若存在 $m\in\R$ 使 $(Su,u)\ge m\|u\|^2$；上半有界可化为 $-S$ 下半有界。
\end{definition}

\begin{theorem}[Friedrichs 扩张]\label{thm:y11-friedrichs}
每个稠密定义、下半有界的对称算子 $S$ 都有保持同一下界的自伴扩张 $S_F$。该扩张由 $S$ 的二次型闭包唯一确定，不意味着所有自伴扩张唯一。
\end{theorem}
\begin{proof}
平移后先设 $S\ge I$，令 $q(u,v)=(Su,v)$，$u,v\in D(S)$。这是内积，且 $\|u\|\le\|u\|_q$。以 $V$ 表示其完备化，映入 $H$ 的连续延拓记为 $j$。关键是 $j$ 单射：若 $u_n$ 在 $q$ 范数中收敛于 $u\in V$，且 $u_n\to0$ 于 $H$，则对 $v\in D(S)$，
\[
 q(u_n,v)=(u_n,Sv)\longrightarrow0.
\]
故 $u$ 与 $V$ 中的稠密子空间 $D(S)$ 正交，$u=0$。这也证明原二次型可闭。因而可把 $V$ 看成 $H$ 中的稠密子空间，并把 $q$ 延为闭形式。

对 $f\in H$，Riesz 表示给出唯一 $Rf\in V$ 使 $q(Rf,v)=(f,v)$。由 $q\ge\|\cdot\|^2$ 得 $R:H\to H$ 有界、正、自伴、单射：单射使用 $V$ 在 $H$ 中稠密；自伴由 $q$ 的对称性得到。故 $R$ 值域稠密。定义 $S_F=R^{-1}$，$D(S_F)=R(H)$。它对称，且 $S_F\pm iI$ 满射，因为 $I\pm iR$ 满射：其范数下界为 $1$，伴随有同一下界，故闭值域且稠密。自伴判据给出 $S_F=S_F^*$。

若 $u\in D(S)$，则 $q(u,v)=(Su,v)$ 先对 $v\in D(S)$ 成立，再延到 $V$，故 $u=RSu$，从而 $S_Fu=Su$。又 $(S_Fu,u)=q(u,u)\ge\|u\|^2$，保持下界。一般 $m$ 情形对 $S+(1-m)I$ 应用此构造再移回。闭形式决定 $R$，故决定此扩张。
\end{proof}

\begin{proposition}[实算子的扩张]
稠密定义的实对称算子有自伴扩张。若还半有界，则 Friedrichs 扩张本身为实算子。
\end{proposition}
\begin{proof}
对闭包仍有实性。共轭将 $\ker(S^*-iI)$ 反线性等距映到 $\ker(S^*+iI)$，故亏指数相等；第 \ref{sec:yosida-11-15} 节的 Cayley 扩张判据给出自伴扩张。半有界时，$C$ 保持型范数和闭型空间 $V$，Riesz 定义的 $R$ 与 $C$ 交换，因此其逆 $S_F$ 与 $C$ 交换。
\end{proof}

\begin{example}[最小 Laplace 算子]
在有界开集 $G$ 上，$S=-\Delta$、$D(S)=C_c^\infty(G)$ 为正对称算子。平移形式 $q_1(u,v)=\int\nabla u\cdot\overline{\nabla v}+\int u\bar v$ 的闭包定义域为 $H_0^1(G)$。其 Friedrichs 扩张满足
\[
 D(S_F)=\{u\in H_0^1(G):-\Delta u\in L^2(G)\},\qquad S_Fu=-\Delta u,
\]
其中导数按分布理解。只有另有边界正则性时，才能再把定义域写成 $H^2(G)\cap H_0^1(G)$。
\end{example}

\begin{exercise}[型范数与图范数]\label{ex:y11-form-graph}
对 $Ae_n=n^2e_n$，比较型域 $D(A^{1/2})$ 与算子域 $D(A)$。
\end{exercise}
\begin{solution}
型域要求 $\sum n^2|x_n|^2<\infty$，算子域要求 $\sum n^4|x_n|^2<\infty$。取 $x_n=n^{-2}$，前者成立、后者不成立。闭型空间不是算子图范数完备化。
\end{solution}

\begin{exercise}[扩张不唯一]\label{ex:y11-friedrichs-nonunique}
在 $L^2(0,1)$ 上解释 $-\dd^2/\dd x^2$ 的最小算子为何同时有 Dirichlet 与 Neumann 自伴扩张，但 Friedrichs 扩张选择前者。
\end{exercise}
\begin{solution}
两种定义域分别为 $H^2\cap H_0^1$ 与 $\{u\in H^2:u'(0)=u'(1)=0\}$；分部积分的边界型和非实边值问题给出自伴性，也可分别用正弦、余弦基对角化。测试函数的能量型闭包是 $H_0^1$，而非 $H^1$，故 Friedrichs 构造选 Dirichlet 扩张。
\end{solution}

\Needspace{10\baselineskip}
\section[自伴算子的谱、Rayleigh 原理和 Krylov--Weinstein 定理、谱的重数]{自伴算子的谱、Rayleigh 原理和\\Krylov--Weinstein 定理、谱的重数}\label{sec:yosida-11-08}
\noindent\textit{原书起始页：中译本 332；英文版 319。}

\begin{theorem}[谱、原子与残差]\label{thm:y11-spectral-tests}
设 $A=A^*$。实数 $\lambda\in\sigma(A)$ 当且仅当每个 $\varepsilon>0$ 都有 $E((\lambda-\varepsilon,\lambda+\varepsilon))\ne0$；并且
\[
 E(\{\lambda\})H=\ker(A-\lambda I),\qquad \sigma_r(A)=\varnothing.
\]
对每个单位向量 $x\in D(A)$ 及实数 $a$，
\[
 \operatorname{dist}(a,\sigma(A))\le\|(A-aI)x\|.
\]
\end{theorem}
\begin{proof}
若某邻域投影为零，$(t-\lambda)^{-1}$ 在测度意义下有界，谱积分给出有界逆；反之，任意小邻域有非零投影时，选其中单位向量 $x$，它属 $D(A)$ 且 $\|(A-\lambda)x\|\le\varepsilon$，与有界逆不相容。原子断言由平方范数积分直接得到。又 $\overline{R(A-\lambda)}^\perp=\ker(A-\lambda)$，故单射时值域必稠密，没有残余谱。最后对 $d=\operatorname{dist}(a,\sigma(A))$，有 $\|(A-a)x\|^2=\int|t-a|^2\,\dd\mu_x\ge d^2$。
\end{proof}

\begin{theorem}[Rayleigh 与 Krylov--Weinstein 估计]\label{thm:y11-rayleigh}
对自伴 $A$，以扩展实数理解，有
\[
 \inf\sigma(A)=\inf_{\|x\|=1,\,x\in D(A)}(Ax,x),\quad
 \sup\sigma(A)=\sup_{\|x\|=1,\,x\in D(A)}(Ax,x).
\]
若 $a=(Ax,x)$、$\|x\|=1$、$r=\|Ax-ax\|$，则 $\sigma(A)$ 中有一点距 $a$ 至多为 $r$；而且
\[
 r^2=\|Ax\|^2-a^2.
\]
\end{theorem}
\begin{proof}
Rayleigh 商是概率测度 $\mu_x$ 的均值，故受谱上下界约束。反向在谱边缘任意小的有界区间取非零投影中的单位向量，其均值就在该区间，逼近谱边缘；无界时取任意远的谱区间，得到扩展实数结论。残差恒等式由展开平方得到，距离界用上一项。实轴上的非空闭谱到固定点的距离可达到，故可把原书带任意正误差的区间估计加强为闭区间估计。
\end{proof}

残差只保证附近有谱，不能无条件保证附近有特征值。只有额外知道该区间的谱离散时，才能把它用于本征值误差估计。

\begin{definition}[简单谱与重数]
若存在 $x$ 使 $\overline{\operatorname{span}\{E(B)x:B\text{ 为 Borel 集}\}}=H$，称 $A$ 有简单谱，$x$ 为循环向量。对于纯点谱，特征值的重数为 $\dim E(\{\lambda\})H$。连续谱处的重数需用循环分解或直积分中的纤维维数描述，不能用单点投影的维数代替。
\end{definition}

例如 $M_t$ 作用于 $L^2(0,1)$ 有简单谱，但所有单点投影均为零；$M_t\oplus M_t$ 在每个频率有两个分量，因而不是简单谱。一般可分情形的测度类与重数函数分类不在本节完整证明，第 \ref{sec:yosida-11-14} 节证明循环情形。

\begin{exercise}[残差不保证特征值]\label{ex:y11-residual}
对 $A=M_t$ 在 $L^2(0,1)$ 上，取归一化指示函数 $x_\varepsilon$ 支于 $(a-\varepsilon,a+\varepsilon)\subset(0,1)$。计算 Rayleigh 商及残差。
\end{exercise}
\begin{solution}
均值为 $a$，方差为 $\varepsilon^2/3$，故残差为 $\varepsilon/\sqrt3$。谱为 $[0,1]$，却没有任何特征值，因为单点测度为零。
\end{solution}

\begin{exercise}[循环与重数二]\label{ex:y11-multiplicity}
证明 $M_t\oplus M_t$ 作用于 $L^2(0,1;\mathbb C^2)$ 不存在循环向量。
\end{exercise}
\begin{solution}
给定 $x=(f,g)$，若 $x=0$ 则显然不循环；否则 $y=(\bar g,-\bar f)$ 非零且属于同一空间。对每个标量 Borel 乘子 $h$，逐点有 $hf\,\overline{\bar g}+hg\,\overline{-\bar f}=0$，故 $y$ 与所有 $h x$ 正交，尤其与全部 $E(B)x$ 正交。
\end{solution}

\Needspace{10\baselineskip}
\section{一般展开定理、连续谱不出现的一个条件}\label{sec:yosida-11-09}
\noindent\textit{原书起始页：中译本 336；英文版 323。}

\begin{theorem}[一般展开与预解式反演]\label{thm:y11-stone-formula}
若 $A=A^*$，则对每个 $x\in H$，$x=\lim_{R\to\infty}E([-R,R])x$。更具体地，对有限 $a<b$，令 $Q(z)=(A-zI)^{-1}$，则
\[
 \lim_{\varepsilon\downarrow0}\frac1{2\pi i}
 \int_a^b\bigl(Q(t+i\varepsilon)-Q(t-i\varepsilon)\bigr)x\,\dd t
 =\bigl(E((a,b))+\tfrac12E(\{a\})+\tfrac12E(\{b\})\bigr)x.
\]
极限为强极限；积分端点如有原子，各计一半。
\end{theorem}
\begin{proof}
对固定 $\varepsilon>0$，有界预解式的 Bochner 积分可与谱积分交换，所得标量乘子为
\[
 k_\varepsilon(s)=\frac1\pi\int_a^b
 \frac{\varepsilon}{(s-t)^2+\varepsilon^2}\,\dd t
 =\frac1\pi\left(\arctan\frac{b-s}{\varepsilon}-\arctan\frac{a-s}{\varepsilon}\right).
\]
它在 $0$ 与 $1$ 之间，逐点趋于开区间指示函数加上两个端点的半质量。用谱积分平方范数公式和控制收敛得到强极限。一般展开来自 $\mu_x([-R,R])\to\mu_x(\R)=\|x\|^2$。
\end{proof}

注意本节 $Q(z)=(A-zI)^{-1}$ 与第八章 $R(z,A)=(zI-A)^{-1}$ 差一个负号；反演公式中的次序也随之固定。

\begin{theorem}[紧自伴展开与紧预解式]\label{thm:y11-compact-expansion}
紧自伴算子 $K$ 的空间可正交分解为其核与所有非零特征空间，且
\[
 x=P_0x+\sum_{\lambda\ne0}P_\lambda x,\qquad
 Kx=\sum_{\lambda\ne0}\lambda P_\lambda x.
\]
后一式按 $|\lambda|>\varepsilon$ 截断时算子范数误差至多 $\varepsilon$。若自伴 $A$ 的一个非实预解式紧，则 $A$ 有完备特征向量系，各特征空间有限维，有界实区间内仅有有限多个特征值；无限特征值序列只能逃向无穷远。
\end{theorem}
\begin{proof}
紧性和第十章说明非零谱离散、各特征空间有限维。其正交和的补空间约化 $K$，且限制算子的谱只能含零；否则非零谱投影会再给出特征向量。自伴性给出 $\|K|_{H_0}\|=r(K|_{H_0})=0$，所以补空间正是核。范数余项由谱乘子在 $|\lambda|\le\varepsilon$ 上至多为 $\varepsilon$ 得到。

令 $Q=(A-iI)^{-1}$，它是紧正规算子且单射。紧正规算子的同一结论由 $\|N\|=r(N)$ 和正交谱投影证明，其非零特征空间正交张成全空间，因为剩余部分范数为零且核为零。$Q$ 的特征值 $\mu\ne0$ 对应 $A$ 的特征值 $i+\mu^{-1}$，且 $Q$ 的特征向量在 $D(A)$ 中。紧性使非零 $\mu$ 只能聚于零，故 $A$ 的特征值只会逃向无穷远。若最初紧的是另一非实预解式，预解恒等式将紧性传到 $i$。
\end{proof}

\begin{example}[有界域上的 Dirichlet Laplace 算子]
第 \ref{sec:yosida-11-07} 节的 $A=-\Delta_D$ 满足 $(A+I)^{-1}:L^2(G)\to H_0^1(G)$ 有界；由 Rellich 紧嵌入，作为 $L^2$ 上算子它紧。因此 $A$ 有正交本征展开。无须在此预先证明 $H^2$ 边界正则性；有界域与零边界型空间已经给出紧性。
\end{example}

\begin{exercise}[端点半质量]\label{ex:y11-stone-endpoints}
若 $A=aI$，求反演公式在区间 $[a,b]$、$b>a$ 上的极限，解释为何不是 $I$。
\end{exercise}
\begin{solution}
谱为集中于 $a$ 的单原子，核乘子在该点趋于 $1/2$，故极限为 $I/2$。要恢复含端点全质量的投影，应另加端点的一半投影，或选避开原子的端点。
\end{solution}

\begin{exercise}[一维本征展开]\label{ex:y11-dirichlet-basis}
对 $(0,\pi)$ 上的 Dirichlet Laplace 算子，写出归一化本征函数、本征值和算子定义域的序列条件。
\end{exercise}
\begin{solution}
$e_n(x)=\sqrt{2/\pi}\sin(nx)$，$\lambda_n=n^2$。写 $u=\sum u_ne_n$，则 $D(A)=\{u:\sum n^4|u_n|^2<\infty\}$，$Au=\sum n^2u_ne_n$。型域只要求 $\sum n^2|u_n|^2<\infty$。
\end{solution}

\Needspace{10\baselineskip}
\section{Peter--Weyl--Neumann 理论}\label{sec:yosida-11-10}
\noindent\textit{原书起始页：中译本 339；英文版 326。}

本节完整证明紧 Hausdorff 群 $G$ 的 Peter--Weyl 定理。采用外部前提：$G$ 有唯一归一化 Haar 概率测度 $m$，它左右不变，且每个非空开集有正测度。原书还对全有界群和抽象群上的殆周期函数使用不变均值；那一推广在此仅作衔接。

\begin{definition}[有限维表示与矩阵系数]
连续酉表示是连续同态 $\pi:G\to U(V_\pi)$，$V_\pi$ 有限维。函数 $g\mapsto(\pi(g)v,w)$ 称矩阵系数，其有限线性组合构成代表函数代数 $\mathcal R(G)$。常数、乘积和共轭分别来自平凡、张量积和共轭表示。
\end{definition}

\begin{lemma}[Schur 正交关系]\label{lem:y11-schur}
每个有限维酉表示是不可约表示的正交直和。取不可约代表 $\pi,\rho$ 的正交基，有
\[
 \int_G\pi_{ij}(g)\overline{\rho_{kl}(g)}\,\dd m(g)
 =\begin{cases}d_\pi^{-1}\delta_{ik}\delta_{jl},&\pi=\rho,\\0,&\pi\not\cong\rho.\end{cases}
\]
\end{lemma}
\begin{proof}
不变子空间的正交补仍不变，按维数归纳得到完全可约性。若 $T$ 是两个不可约表示间的交织算子，其核与像不变，所以非零时为同构；同一不可约表示的交织算子取一个复特征值，$T-\lambda I$ 的核非零，故 $T=\lambda I$。这就是 Schur 引理。

给定线性映射 $B:V_\rho\to V_\pi$，平均 $\int\pi(g)B\rho(g)^{-1}\,\dd m(g)$ 为交织算子。不同不可约时它为零；相同时它为 $(\operatorname{tr}B/d_\pi)I$，因取迹不改变。对矩阵单位 $B$ 比较各矩阵元，便得公式。
\end{proof}

\begin{theorem}[Peter--Weyl]\label{thm:y11-peter-weyl}
$\mathcal R(G)$ 在 $C(G)$ 中一致稠密；各不可约表示的归一化系数 $\sqrt{d_\pi}\,\pi_{ij}$ 构成 $L^2(G,m)$ 的完备正交系。尤其有限维连续酉表示分离 $G$ 中的点。
\end{theorem}
\begin{proof}
先构造紧的平滑算子。取连续实函数 $k\ge0$，$k(g^{-1})=k(g)$，支于单位元的对称小邻域，且 $\int k=1$。定义
\[
 T_kf(x)=\int_G k(y^{-1}x)f(y)\,\dd m(y).
\]
核连续，故 $T_k:L^2(G)\to C(G)$ 有界：Cauchy--Schwarz 给出一致界，核的平移在 $L^2$ 中连续给出连续性。它在 $L^2$ 上为紧自伴算子，并与左平移 $L_af(x)=f(a^{-1}x)$ 交换。所有非零特征空间有限维、左平移不变，其中的函数连续。因此左平移在这些空间上给出有限维连续酉表示，而其中每个函数通过在单位元处取值可写成该表示或其对偶表示的矩阵系数。

若 $v\in L^2$ 与所有有限维系数正交，紧自伴展开使 $T_kv=0$。当单位元邻域缩小时，$T_kv\to v$ 于 $L^2$：先对连续函数用一致连续性，再由它们在 $L^2$ 中稠密及 $\|T_k\|\le1$ 推广。故 $v=0$，证明矩阵系数 $L^2$ 稠密。结合 Schur 关系即得完备正交系。

最后取 $f\in C(G)$，先选 $k$ 使 $\|T_kf-f\|_\infty<\varepsilon$，再选 $p\in\mathcal R(G)$ 使 $\|f-p\|_2<\varepsilon/\|T_k\|_{2\to\infty}$。卷积一个矩阵系数仍为同一有限维表示的系数线性组合，所以 $T_kp\in\mathcal R(G)$，且 $\|T_kp-f\|_\infty<2\varepsilon$。这证明一致稠密。若表示不分离两个点，所有代表函数均在两点取相同值，与连续函数分离点及一致稠密矛盾。
\end{proof}

不可分 $L^2(G)$ 的正交基可不可数，展开按有限子集的网理解。圆群结论只保证三角多项式一致稠密，不保证普通 Fourier 部分和一致收敛。

\begin{exercise}[紧交换群的不可约表示]\label{ex:y11-abelian-representations}
证明紧交换群的有限维不可约酉表示都是一维；对圆群写出这些表示。
\end{exercise}
\begin{solution}
交换性使每个 $\pi(g)$ 都是交织算子，Schur 引理使它为标量。若空间维数大于一，任意一维子空间都不变，违反不可约性。圆群的一维连续特征为 $e^{i\theta}\mapsto e^{in\theta}$，$n\in\mathbb Z$，由实轴特征及 $2\pi$ 周期性得到。
\end{solution}

\begin{exercise}[有限群的维数公式]\label{ex:y11-finite-group}
若 $G$ 有限，证明 $\sum_{\pi\in\widehat G}d_\pi^2=|G|$。
\end{exercise}
\begin{solution}
归一化计数测度下 $L^2(G)$ 的维数为 $|G|$；Peter--Weyl 的正交基每个不可约表示贡献 $d_\pi^2$ 个系数，数基向量便得等式。
\end{solution}

\Needspace{10\baselineskip}
\section{关于非交换紧群的 Tannaka 对偶性定理}\label{sec:yosida-11-11}
\noindent\textit{原书起始页：中译本 345；英文版 332。}

交换紧群可用一维特征描述，非交换紧群则需要所有有限维表示及其乘法关系。下面采用与原书同型的代表函数环表述，避免把“任意一组酉矩阵”误认为群的一个点。

\begin{theorem}[Tannaka 的代表函数表述]\label{thm:y11-tannaka}
设 $G$ 为紧 Hausdorff 群，$\mathcal R(G)$ 为代表函数代数。每个含单位、保持复共轭的复代数同态
\[
 \chi:\mathcal R(G)\longrightarrow\mathbb C
\]
都是唯一的点值泛函 $\chi(f)=f(g)$。在所有此类 $\chi$ 上赋逐点收敛拓扑及下述群运算，点值映射给出 $G$ 到该空间的拓扑群同构。这里没有预先假定 $\chi$ 对一致范数连续。
\end{theorem}
\begin{proof}
对每个不可约酉表示 $\pi$，矩阵 $U_\pi(\chi)=(\chi(\pi_{ij}))$ 是酉矩阵：直接把 $\pi(g)^*\pi(g)=I$ 的各系数恒等式送入保持共轭的 $\chi$。代表函数由这些系数线性张成，所以 $\chi$ 由所有 $U_\pi$ 决定。线性关系、乘法关系、共轭及单位关系均为闭条件，故所有这样的 $\chi$ 组成乘积 $\prod_\pi U(d_\pi)$ 的闭子集 $K$，从而紧。

代表函数有代数余乘法 $\Delta f(g,h)=f(gh)$，其值在 $\mathcal R(G)\otimes\mathcal R(G)$ 中，因为 $\Delta\pi_{ij}=\sum_k\pi_{ik}\otimes\pi_{kj}$。定义
\[
 (\chi\psi)(f)=(\chi\otimes\psi)(\Delta f),\qquad
 \chi^{-1}(f)=\chi(f\circ\mathrm{inv}).
\]
这些公式在系数上就是矩阵乘法与取逆，代数恒等式保证它们仍为含单位、保持共轭的同态。于是 $K$ 为紧拓扑群。点值映射 $G\to K$ 连续，Peter--Weyl 分离点使它单射；紧性使其像为闭子群。

将 $f\in\mathcal R(G)$ 变成 $\widetilde f(\chi)=\chi(f)$。这些函数在 $K$ 上含常数、分离点且对共轭封闭，所以一致稠密于 $C(K)$。每个 $\pi$ 的矩阵 $U_\pi$ 是 $K$ 上的连续酉表示；它限制到 $G$ 后是不可约 $\pi$，所以在 $K$ 上也不可约，且非平凡表示仍非平凡。对 $K$ 和 $G$ 分别使用 Haar 平均：平凡表示系数的平均为 $1$，其他不可约系数的平均均为零。因此两种平均对每个 $\widetilde f$ 相同，再由一致稠密性对所有 $C(K)$ 函数相同。

若 $G$ 的像是真闭子群，Urysohn 引理给出非负连续函数在其上为零、在像外某点为正。它对 $G$ 的积分为零，对 $K$ 的积分为正，因为非空开集 Haar 测度为正，矛盾。故 $G=K$，得到满射；唯一性及拓扑同构随分离点和紧 Hausdorff 性得到。
\end{proof}

等价地，群的一个点作用在每个有限维表示上，且与所有交织算子、直和、张量积、共轭及平凡表示相容。代表函数中的代数关系正编码这些条件。仅要求每个表示对应一个酉矩阵，没有这些相容条件，会得到远大于 $G$ 的乘积空间。

\begin{exercise}[圆群的重构]\label{ex:y11-tannaka-circle}
对 $G=\mathbb T$，用 Laurent 多项式代数证明上述重构。
\end{exercise}
\begin{solution}
$\mathcal R(\mathbb T)=\mathbb C[z,z^{-1}]$，共轭满足 $z^*=z^{-1}$。同态由 $w=\chi(z)$ 决定，保持共轭及 $zz^{-1}=1$ 给出 $|w|=1$，于是 $\chi(p)=p(w)$，恰是点值。
\end{solution}

\begin{exercise}[张量相容性]\label{ex:y11-tannaka-tensor}
圆群的第 $n$ 个特征记为 $\pi_n$。若为每个 $n$ 任取 $u_n\in\mathbb T$，写出来自同一群元素所需的条件。
\end{exercise}
\begin{solution}
需 $u_0=1$、$u_{m+n}=u_mu_n$、$u_{-n}=\overline{u_n}$。这些条件使 $u_n=u_1^n$，于是由唯一的 $u_1\in\mathbb T$ 决定。独立选择各个 $u_n$ 通常不满足乘法关系。
\end{solution}

\Needspace{20\baselineskip}
\section{自伴算子的函数}\label{sec:yosida-11-12}
\noindent\textit{原书起始页：中译本 351；英文版 338。}

令 $A=\int t\,\dd E(t)$。对有限值 Borel 函数 $f$，定义 $f(A)$ 为定理 \ref{thm:y11-spectral-integral} 中的闭算子。允许在 $E$-零集上修改函数。其定义域取决于向量的谱测度，而不仅取决于形式表达式。

\begin{theorem}[无界函数演算的运算规则]\label{thm:y11-borel-calculus}
设 $f,g$ 为有限值 Borel 函数。则 $f(A)^*=\bar f(A)$，并且
\begin{align*}
 D(g(A)f(A))&=D(f(A))\cap D((gf)(A)),\\
 g(A)f(A)x&=(gf)(A)x\quad(x\text{ 属于左侧定义域}),\\
 (f+g)(A)x&=f(A)x+g(A)x\quad(x\in D(f(A))\cap D(g(A))).
\end{align*}
$f(A)$ 正规；$f$ 实值时自伴。复合满足 $h(f(A))=(h\circ f)(A)$，其中复值 $f$ 时左侧用正规算子的 Borel 演算，并取对应的最大定义域。
\end{theorem}
\begin{proof}
对谱投影 $E(B)$ 和 $x\in D(f(A))$，截断后取极限得 $E(B)f(A)x=f(A)E(B)x$，且
\[
 \mu_{f(A)x}(B)=\int_B|f|^2\,\dd\mu_x.
\]
故 $f(A)x\in D(g(A))$ 等价于 $\int|gf|^2\,\dd\mu_x<\infty$，连同先验的 $x\in D(f(A))$ 得出定义域公式。等式对共同截断成立，再由平方积分极限推广。和的公式同理，但一般只能在定义域交上断言。

伴随等式已证。$f(A)^*f(A)$ 与 $f(A)f(A)^*$ 的定义域均为 $D(|f|^2(A))$，因为四次可积结合 $x\in H$ 蕴含二次可积；两者都等于 $|f|^2(A)$，故正规。$f(A)$ 的谱测度是 $E$ 经 $f$ 的推前，复合的定义域条件就是 $\int|h(f(t))|^2\,\dd\mu_x<\infty$，由截断给出复合公式。
\end{proof}

例如若 $f(t)=t$、$g(t)=1/t$（在零处另行有限定义），即使乘积在谱上等于 $1$，$g(A)f(A)$ 仍只在 $D(A)$ 上定义，未必是处处定义的恒等算子。它的闭包可以是 $I$，而算子本身不等于 $I$。

\begin{proposition}[与谱投影交换]\label{prop:y11-commutation}
有界算子 $B$ 若与自伴 $A$ 在定义域上交换，即 $BD(A)\subset D(A)$ 且 $ABx=BAx$，则与每个 $E(C)$ 交换；因此保持 $D(f(A))$ 并满足 $Bf(A)x=f(A)Bx$。
\end{proposition}
\begin{proof}
定义域上的关系先给出 $B(A\pm iI)^{-1}=(A\pm iI)^{-1}B$，从而与 Cayley 变换及其伴随交换。连续演算、再以有界 Borel 函数的强逼近得到 $BE(C)=E(C)B$。为处理无界 $f$，令 $P_n=E(\{|f|\le n\})$，则 $f(A)P_nBx=Bf(A)P_nx$，右侧收敛，左侧由闭性得到 $Bx\in D(f(A))$ 及交换等式。
\end{proof}

原书进一步证明 Neumann--Riesz--Mimura 的闭算子判别：在可分空间上，适当的交换子包含关系刻画 $T=f(A)$。这里证明正向的谱交换关系及循环模型中的乘法解释；一般闭算子判别的反向拼接证明尚未展开，不能把“与 $A$ 交换”本身当作“是 $A$ 的标量函数”的充分条件，重数大于一时尤其如此。

\begin{exercise}[定义域使恒等式失效]\label{ex:y11-product-domain}
在 $L^2(1,\infty)$ 上取 $A=M_t$。比较 $A^{-1}A$ 与 $AA^{-1}$ 的定义域。
\end{exercise}
\begin{solution}
$A^{-1}=M_{1/t}$ 有界且值域为 $D(A)$。所以 $AA^{-1}=I$ 处处定义；$A^{-1}A=I|_{D(A)}$，只在 $D(A)$ 定义。其闭包才是 $I$。
\end{solution}

\begin{exercise}[实函数与复函数]\label{ex:y11-functional-examples}
对 $A\ge0$ 定义 $A^{1/2}$ 与 $e^{-tA}$，证明前者自伴，后者在 $t\ge0$ 时为压缩算子且强连续。
\end{exercise}
\begin{solution}
使用乘子 $\sqrt\lambda$ 与 $e^{-t\lambda}$。前者实值，定义域要求 $\int\lambda\,\dd\mu_x<\infty$，故自伴。后者绝对值至多 $1$，给出压缩；$t\to s\ge0$ 时乘子逐点收敛且有界，控制收敛给出强连续。
\end{solution}

\Needspace{10\baselineskip}
\section{Stone 定理和 Bochner 定理}\label{sec:yosida-11-13}
\noindent\textit{原书起始页：中译本 358；英文版 345。}

\begin{theorem}[Stone 定理的谱积分形式]\label{thm:y11-stone}
若 $A=A^*$，则
\[
 U(t)=e^{itA}=\int_\R e^{it\lambda}\,\dd E(\lambda)
\]
为强连续酉群，生成元为 $iA$，包括定义域相等。每个强连续酉群都唯一以此形式表示。
\end{theorem}
\begin{proof}
乘子乘法和 $|e^{it\lambda}|=1$ 给出群律与酉性。控制收敛给出强连续。对 $x\in D(A)$，$|(e^{it\lambda}-1)/t|\le|\lambda|$，差商在平方积分中趋于 $i\lambda$，故生成元包含 $iA$。

反之，若差商在 $H$ 中收敛，则其范数有界。沿 $t_n\to0$ 用 Fatou 引理，得到 $\int\lambda^2\,\dd\mu_x<\infty$，所以 $x\in D(A)$，生成元正好为 $iA$。对任意强连续酉群，第九章定理 \ref{thm:y9-stone} 已在不使用谱定理的前提下证明生成元为某个 $iA$。本节积分群与它生成元相同，由生成元唯一性两群相同。
\end{proof}

\begin{definition}[正定函数]
函数 $\varphi:\R\to\mathbb C$ 称为正定，若对任意有限组 $t_j\in\R$、$c_j\in\mathbb C$，
\[
 \sum_{j,k}c_j\overline{c_k}\varphi(t_j-t_k)\ge0.
\]
该条件蕴含 $\varphi(0)\ge0$、$\varphi(-t)=\overline{\varphi(t)}$ 及 $|\varphi(t)|\le\varphi(0)$，可由一阶与二阶矩阵检验得到。
\end{definition}

\begin{theorem}[Bochner 表示]\label{thm:y11-bochner}
$\varphi$ 是连续正定函数，当且仅当存在唯一有限正 Borel 测度 $\mu$ 使
\[
 \varphi(t)=\int_\R e^{it\lambda}\,\dd\mu(\lambda),\qquad
 \mu(\R)=\varphi(0).
\]
\end{theorem}
\begin{proof}
若存在测度，正定性由积分 $|\sum_jc_je^{it_j\lambda}|^2$ 得到，连续性由控制收敛得到。

反向在有限形式和 $\sum c_j\delta_{t_j}$ 上定义半内积 $(\delta_s,\delta_t)=\varphi(s-t)$。正定性保证非负；Cauchy--Schwarz 使零范数向量构成可商掉的子空间。商后完备化为 Hilbert 空间 $K$，平移 $V(a)\delta_s=\delta_{s+a}$ 保内积，逆为 $V(-a)$。又
\[
 \|\delta_{s+a}-\delta_s\|^2=2\varphi(0)-2\operatorname{Re}\varphi(a)\longrightarrow0,
\]
故先在稠密有限和上、再由等距性在全空间上得到强连续酉群。记 $v=\delta_0$ 的商像，则 $\varphi(t)=(V(t)v,v)$。对 $V(t)$ 使用 Stone 谱积分，取 $\mu(B)=(E(B)v,v)$ 即得表示。

若两有限测度有同一 Fourier 变换，对任意 Schwartz 测试函数应用 Fourier 反演与 Fubini，得到它们对所有 $C_c^\infty$ 函数的积分相同；有限测度的正则性使它们相等。因此表示唯一。
\end{proof}

这里正定与逐点非负不同；概率测度对应额外条件 $\varphi(0)=1$。积分指数取正号，与本书函数 Fourier 变换的负号约定配套。

\begin{exercise}[含负值的正定函数]\label{ex:y11-positive-definite}
证明 $\cos t$ 正定，求其表示测度，并说明正定不意味着逐点非负。
\end{exercise}
\begin{solution}
$\cos t=(e^{it}+e^{-it})/2$，对应正测度 $(\delta_1+\delta_{-1})/2$，所以正定；但 $\cos\pi=-1$。
\end{solution}

\begin{exercise}[有界生成元与范数连续]\label{ex:y11-unitary-norm-continuity}
证明 $e^{itA}$ 在 $t=0$ 范数连续当且仅当自伴 $A$ 有界。
\end{exercise}
\begin{solution}
有界时指数级数给出范数连续。若 $A$ 无界，谱中存在 $|\lambda_n|\to\infty$，取 $t_n=\pi/|\lambda_n|$，则 $|e^{it_n\lambda_n}-1|=2$。连续演算的等距性或谱支集判别给出 $\|e^{it_nA}-I\|=2$，故不范数连续。
\end{solution}

\Needspace{10\baselineskip}
\section{具有简单谱的自伴算子的标准型}\label{sec:yosida-11-14}
\noindent\textit{原书起始页：中译本 360；英文版 347。}

\begin{theorem}[循环谱的乘法模型]\label{thm:y11-cyclic-model}
若自伴 $A$ 有循环向量 $v$，令 $\mu(B)=(E(B)v,v)$。则
\[
 W:L^2(\R,\mu)\longrightarrow H,\qquad Wf=f(A)v
\]
先对有界简单函数定义，再延拓为酉映射，并满足
\[
 W^{-1}AW=M_t,\qquad
 D(M_t)=\{f\in L^2(\mu):tf\in L^2(\mu)\}.
\]
谱投影变成乘 $1_B$。此处 $v$ 不必属于 $D(A)$。
\end{theorem}
\begin{proof}
对简单函数 $f,g$，正交投影关系给出 $(Wf,Wg)=\int f\bar g\,\dd\mu$。简单函数在 $L^2(\mu)$ 中稠密，故可等距延拓；其像闭，且包含所有 $E(B)v$，循环性使其满射。又 $E(B)Wf=W(1_Bf)$，从而 $Wf$ 的谱测度为 $|f|^2\mu$。因此 $Wf\in D(A)$ 恰等价于 $tf\in L^2(\mu)$；对谱截断先验证 $AWf=W(tf)$，再由闭性得全定义域上的等式。
\end{proof}

\begin{proposition}[简单谱的酉不变量]\label{prop:y11-measure-class}
两个有限正测度的坐标乘法算子 $M_t$ 具有循环向量 $1$，且它们酉等价当且仅当这两个测度相互绝对连续。因而谱集相同不足以确定自伴算子的酉等价类。
\end{proposition}
\begin{proof}
常数 $1$ 循环，因为简单函数稠密。若 $\mu\sim\nu$，以 $f\mapsto\sqrt{\dd\mu/\dd\nu}\,f$ 定义酉映射；它与乘 $t$ 交换，并保持对应定义域。反之，若 $U:L^2(\mu)\to L^2(\nu)$ 交织两个自伴乘法算子，它也交织谱投影。置 $h=U1$，则 $\mu(B)=\|U1_B\|^2=\int_B|h|^2\,\dd\nu$。因为 $h$ 循环，它不能在一个正 $\nu$ 测度集合上为零，故 $\mu\sim\nu$。
\end{proof}

例如 $[0,1]$ 上的 Lebesgue 测度与 $\nu=\sum_n2^{-n}\delta_{q_n}$（$q_n$ 枚举互异有理数）均有全支集 $[0,1]$。两个乘法算子的谱都为 $[0,1]$，但前者无点谱、后者为纯点谱，不酉等价。

\begin{exercise}[换循环向量]\label{ex:y11-change-cyclic}
在 $L^2(\mu)$ 的坐标乘法模型中，证明 $h$ 为循环向量当且仅当 $h\ne0$ 几乎处处。
\end{exercise}
\begin{solution}
若 $h$ 在正测度集上为零，所有 $1_Bh$ 也为零，无法稠密。反之，对任意 $f\in L^2(\mu)$，令 $f_n=f1_{\{|f|\le n,\,1/n\le|h|\le n\}}$。则 $f_n/h$ 有界，故 $f_n$ 属于有界 Borel 乘子作用于 $h$ 的闭包；$f_n\to f$，故循环。
\end{solution}

\begin{exercise}[离散简单谱]\label{ex:y11-discrete-cyclic}
对 $Ae_n=ne_n$ 作用于 $\ell^2$，给出一个循环向量及对应有限谱测度。
\end{exercise}
\begin{solution}
取 $v=\sum_{n\ge1}2^{-n}e_n$。每个 $E(\{n\})v=2^{-n}e_n$，故循环；对应 $\mu=\sum_{n\ge1}4^{-n}\delta_n$。此例还说明循环向量不要求范数为 $1$；测度总质量为 $\|v\|^2$。
\end{solution}

\Needspace{10\baselineskip}
\section{对称算子的亏指数、广义单位分解}\label{sec:yosida-11-15}
\noindent\textit{原书起始页：中译本 362；英文版 349。}

设 $S$ 闭、稠密定义且对称，沿用第七章
\[
 N_+=\ker(S^*-iI),\quad N_-=\ker(S^*+iI),\quad
 n_\pm=\dim N_\pm.
\]
其 Cayley 变换从 $M_+=R(S+iI)$ 等距映到 $M_-=R(S-iI)$，且 $M_\pm^\perp=N_\pm$。亏指数是 Hilbert 维数，允许无穷；取相反 Cayley 约定会交换两者，必须同时改动符号。

\begin{theorem}[自伴扩张与最大对称性]\label{thm:y11-extensions}
$S$ 在同一 $H$ 中有自伴扩张当且仅当 $n_+=n_-$。其自伴扩张与所有酉映射 $V:N_+\to N_-$ 一一对应，定义为
\begin{align*}
 D(S_V)&=D(S)\mathbin{\dot+}\{h+Vh:h\in N_+\},\\
 S_V(u+h+Vh)&=Su+ih-iVh.
\end{align*}
$S$ 为最大对称算子当且仅当 $n_+=0$ 或 $n_-=0$；最大对称未必自伴。$S$ 自伴当且仅当两亏指数皆为零。
\end{theorem}
\begin{proof}
任何对称扩张对应 Cayley 等距映射的扩张。它为自伴恰对应定义域与值域均为 $H$ 的酉扩张，所以须以酉映射配对缺少的两个正交补；这恰等价于亏指数相等。任一补空间的酉配对都给出酉扩张 $\widetilde U$。因为 $R(I-\widetilde U)$ 包含稠密的 $D(S)$，$\widetilde U$ 不含不动向量，逆 Cayley 公式有效。将补空间配对取为 $-V$，代入 $(I-\widetilde U)h=h+Vh$、$i(I+\widetilde U)h=ih-iVh$，得到所写定义域及作用式。反向每个自伴扩张都如此取得。

若两个缺陷空间均非零，可在各取一维等距配对，扩张原 Cayley 映射，得到严格的对称扩张。若其中一个为零，则原等距映射的定义域或值域已经为全空间，无法再扩大定义域而保持等距性，故最大对称。最后一项就是第七章自伴判据。
\end{proof}

\begin{theorem}[扩空间与广义谱测度]\label{thm:y11-generalized-resolution}
每个闭稠密对称算子 $S$ 都有某个更大 Hilbert 空间 $\widetilde H\supset H$ 中的自伴扩张 $\widetilde A$。若 $P_H$ 为正交投影，则
\[
 F(B)=P_HE_{\widetilde A}(B)|_H
\]
是正算子值测度，$F(\R)=I_H$，并且对 $x\in D(S)$，
\[
 Sx=\int t\,\dd F(t)x,\qquad
 \int t^2\,\dd(F(t)x,x)=\|Sx\|^2.
\]
一般 $F(B)$ 不是投影；二阶积分有限可能只保证 $x\in D(\widetilde A)\cap H$，未必 $x\in D(S)$。若 $S$ 最大对称，则该条件等价于 $x\in D(S)$。
\end{theorem}
\begin{proof}
取 $\widetilde H=H\oplus H$，在其上考虑 $S\oplus(-S)$。亏指数均为 $n_++n_-$，故有自伴扩张 $\widetilde A$。把原 $H$ 视为第一分量便得到所需扩张。压缩谱测度保持正性、单位质量和强可加性。对 $x\in D(S)$，$\widetilde Ax=Sx\in H$，所以压缩谱积分给出作用式，二阶矩等于 $\|\widetilde Ax\|^2$。

对于任意 $x\in H$，二阶矩有限等价于 $x\in D(\widetilde A)\cap H$。算子 $T=P_H\widetilde A|_{D(\widetilde A)\cap H}$ 对称且扩张 $S$；若 $S$ 最大对称，则 $T=S$，从而定义域相等。一般情形没有这个最大性推论。
\end{proof}

\begin{example}[区间与半直线上的动量]
对 $S=-i\,\dd/\dd x$，$D(S)=H_0^1(0,1)$，伴随定义域为 $H^1(0,1)$。解 $S^*u=\pm iu$ 得 $e^{-x}$ 与 $e^x$，均平方可积，所以亏指数为 $(1,1)$。自伴扩张可写成边界条件 $u(1)=e^{i\theta}u(0)$，本征值为 $\theta+2\pi n$。

在半直线 $(0,\infty)$ 上，$e^{-x}$ 可积而 $e^x$ 不可积，故同一符号约定下亏指数为 $(1,0)$，算子最大对称但无同空间自伴扩张。可把它扩到整直线的自伴动量算子，再压缩其谱测度。
\end{example}

\begin{exercise}[压缩未必为投影]\label{ex:y11-compressed-pvm}
在 $\mathbb C^2$ 上取 $\widetilde A=\operatorname{diag}(-1,1)$，把 $H$ 取为 $(1,1)$ 张成的直线。求 $F(\{1\})$ 并检验幂等性。
\end{exercise}
\begin{solution}
以 $(1,1)/\sqrt2$ 为 $H$ 的单位基，$F(\{1\})=\tfrac12I_H$，其平方为 $\tfrac14I_H$，不是投影。这说明广义单位分解应理解为正算子值测度。
\end{solution}

\begin{exercise}[区间边界相位]\label{ex:y11-boundary-phase}
对 $u(1)=e^{i\theta}u(0)$ 的动量算子写出归一化本征函数，并解释不同 $\theta$ 给出不同扩张。
\end{exercise}
\begin{solution}
$u_n(x)=e^{i(\theta+2\pi n)x}$，范数为 $1$。乘 $e^{i\theta x}$ 将周期 Fourier 基变成该基，故完备。不同模 $2\pi$ 的相位给出不同边界定义域及不同本征值序列；它们都包含 $H_0^1$ 上的最小算子。
\end{solution}

\Needspace{10\baselineskip}
\section{群环 \texorpdfstring{$L^1$}{L1} 及 Wiener 的 Tauber 定理}\label{sec:yosida-11-16}
\noindent\textit{原书起始页：中译本 367；英文版 354。}

在 $L^1(\R)$ 上取卷积乘法和 $f^*(x)=\overline{f(-x)}$。Fubini 给出 $\|f*g\|_1\le\|f\|_1\|g\|_1$，所以这是交换 Banach 代数，但没有卷积单位：形式上的单位 $\delta_0$ 不属于 $L^1$。以下用 $\widetilde{L^1}=\mathbb C1\oplus L^1$ 表示单位化。

\begin{theorem}[群环的特征与半单性]\label{thm:y11-group-characters}
$L^1(\R)$ 的所有非零特征为
\[
 \chi_\xi(f)=\widehat f(\xi)=\int_\R e^{-it\xi}f(t)\,\dd t,\qquad \xi\in\R.
\]
单位化的特征另外包含 $\chi_\infty(\lambda1+f)=\lambda$，其他特征为 $\lambda+\widehat f(\xi)$。该代数半单。
\end{theorem}
\begin{proof}
非零特征延为单位化的含单位特征，故自动连续。选 $g$ 使 $\chi(g)\ne0$，令 $\tau_tg(x)=g(x-t)$，$c(t)=\chi(\tau_tg)/\chi(g)$。恒等式 $(\tau_tg)*h=g*(\tau_th)$ 说明此比值与所选非零特征值函数无关，且 $\chi(\tau_th)=c(t)\chi(h)$ 对所有 $h$ 成立。于是 $c(t+s)=c(t)c(s)$，$c(0)=1$。平移的 $L^1$ 连续性给出 $c$ 连续，等距性给出全轴有界。群律使 $|c(t)|=1$，故 $c(t)=e^{-it\xi}$，其中实数 $\xi$ 唯一。

Bochner 积分恒等式 $f*g=\int f(t)\tau_tg\,\dd t$ 使
\[
 \chi(f)\chi(g)=\chi(g)\int f(t)c(t)\,\dd t.
\]
消去 $\chi(g)$ 得 Fourier 公式。反向由 Fubini 验证这些都是特征；单位化中限制到 $L^1$ 为零时只剩 $\chi_\infty$。若所有特征在 $\lambda1+f$ 上为零，则先由 $\chi_\infty$ 得 $\lambda=0$，再由 Fourier 变换的单射性得 $f=0$，故根为零。
\end{proof}

\begin{theorem}[Wiener 平移稠密定理]\label{thm:y11-wiener}
对 $f\in L^1(\R)$，其所有平移的线性张成在 $L^1(\R)$ 中稠密，当且仅当 $\widehat f(\xi)\ne0$ 对所有实数 $\xi$ 成立。
\end{theorem}
\begin{proof}
必要性：若 $\widehat f(\xi_0)=0$，连续非零泛函 $\chi_{\xi_0}$ 消灭所有平移及其闭包，故不稠密。

充分性：先证明闭卷积理想 $I=\overline{f*L^1}$ 为全空间。取 $\varphi$ 满足 $\widehat\varphi\in C_c^\infty(\R)$。选 $h$ 使 $\widehat h\in C_c^\infty$、$0\le\widehat h\le1$，且在 $\operatorname{supp}\widehat\varphi$ 上恒为 $1$；由 Fourier 反演 $h\in\Sch\subset L^1$。
\par\Needspace{4\baselineskip}
在单位化代数中令
\[
 d=1-h+f*f^*.
\]
每个有限特征给出 $1-\widehat h(\xi)+|\widehat f(\xi)|^2>0$，而无穷特征值为 $1$。极大理想判据使 $d$ 可逆。置 $g=d^{-1}*(f^* *\varphi)\in L^1$。在 $\widehat\varphi$ 的支集上 $\widehat h=1$，所以 $\widehat{f*g}=\widehat\varphi$；支集外两者均为零。Fourier 单射性给出 $f*g=\varphi$，故所有这类 $\varphi$ 均在 $I$ 中。

这些 $\varphi$ 在 $L^1$ 中稠密。具体地，先把任意 $L^1$ 函数逼近为 Schwartz 函数 $v$，再取 $\eta\in\Sch$ 满足 $\widehat\eta\in C_c^\infty$、$\widehat\eta(0)=1$。$\eta_\varepsilon(x)=\varepsilon^{-1}\eta(x/\varepsilon)$ 为有界 $L^1$ 近似单位，故 $v*\eta_\varepsilon\to v$ 于 $L^1$，而其 Fourier 变换光滑且紧支。因此 $I=L^1$。

最后 $f*g$ 是平移 $f$ 的 Bochner 加权积分，所以属于平移张成的闭包；可由简单函数逼近积分验证。反向每个 $\tau_af$ 都由 $f*(\tau_a\eta_\varepsilon)$ 逼近。因此闭理想恰为平移张成的闭包，结论成立。
\end{proof}

这里使用的是 $f$ 的平移稠密性，不是 $f$ 在 $L^1$ 卷积代数中可逆。事实上 $\widehat f(\xi)\to0$ 于无穷远，不能把它的倒数当作某个 $L^1$ 函数的 Fourier 变换。

\begin{corollary}[有界函数的 Tauber 传递]\label{cor:y11-tauber-transfer}
若 $h\in L^\infty(\R)$、$\widehat f$ 处处非零，并且
\[
 (f*h)(x)\longrightarrow C\int f\quad(x\to+\infty),
\]
则对每个 $g\in L^1$，$(g*h)(x)\to C\int g$。
\end{corollary}
\begin{proof}
结论对 $f$ 的平移成立，因为 $(\tau_af*h)(x)=(f*h)(x-a)$，故对它们的有限线性组合成立。用 Wiener 稠密定理逼近 $g$，并估计 $\|(g-p)*h\|_\infty\le\|g-p\|_1\|h\|_\infty$、$|\int(g-p)|\le\|g-p\|_1$，即可先控制逼近误差，再令 $x\to\infty$。
\end{proof}

原书末尾还讨论 Littlewood 的 Abel--Tauber 推论。本节未展开其证明；通常级数收敛还需 $na_n$ 有界等条件，不能仅凭 Abel 可和性推出。

\begin{exercise}[高斯平移]\label{ex:y11-gaussian-translates}
证明 $f(x)=e^{-x^2}$ 的有限平移线性组合在 $L^1(\R)$ 中稠密。
\end{exercise}
\begin{solution}
由本书 Fourier 约定，$\widehat f(\xi)=\sqrt\pi e^{-\xi^2/4}$，处处非零，因此应用 Wiener 定理。结论允许复系数的有限线性组合，不要求正系数。
\end{solution}

\begin{exercise}[一个零点就足够阻碍稠密]\label{ex:y11-zero-obstruction}
取 $f=1_{[0,1]}$。证明其平移张成不在 $L^1(\R)$ 中稠密，并写出一个消灭该闭包的非零连续泛函。
\end{exercise}
\begin{solution}
$\widehat f(2\pi)=\int_0^1e^{-2\pi ix}\,\dd x=0$。泛函 $g\mapsto\int e^{-2\pi ix}g(x)\,\dd x$ 非零且有界，并在每个平移 $f$ 上取零，故其核包含整个闭包。
\end{solution}


\Needspace{6\baselineskip}
\section*{本章参考说明与阅读衔接}
\phantomsection\addcontentsline{toc}{section}{本章参考说明与阅读衔接}
直接来源为吉田第十一章 \cite{yosida-zh,yosida-en}。本章采用循环子空间构造谱测度，以紧群卷积和 Fourier 局部逆组织后半章证明；与原书的范畴集、殆周期均值及赋范环路线作适用范围对照。原书所引 Gelfand、Neumann、Friedrichs、Tannaka、Wiener 等原始文献未计为直接阅读。闭算子的可测函数判别一般版本、非可分谱重数分类及 Littlewood 的最终 Tauber 推论等未展开范围见 COVERAGE。下一章转向凸集端点、向量格及其表示。
